\documentclass[
    a4paper,
    man,
    donotrepeattitle,
    floatsintext,
    british
]{apa6}

\usepackage{standalone}
\usepackage{newclude}

\usepackage[british]{babel}
\usepackage[utf8]{inputenc}
\usepackage{epstopdf}
\usepackage{csquotes}
\usepackage[hidelinks]{hyperref}
\usepackage[
    style=apa,
    backend=biber,
    sortcites=true,
    sorting=nyt,
%    isbn=false,
%    url=false,
%    doi=false,
%    eprint=false,
%    hyperref=false,
%    backref=false,
%    firstinits=false,
]{biblatex}

\DeclareLanguageMapping{british}{british-apa}

% maps apacite commands to biblatex commands
\let \citeNP \cite
\let \citeA \textcite
\let \cite \parencite

\bibliography{references}

\usepackage{amsmath}
\usepackage{graphicx}
\usepackage{hyperref}
\usepackage{subcaption}
\hypersetup{colorlinks,urlcolor=blue}

\usepackage{listings}
\usepackage{xcolor}
\definecolor{darkgray}{rgb}{.35,.25,.35}
\definecolor{lightgray}{rgb}{.6,.6,.6}

\lstdefinelanguage{bash}{
  basicstyle=\scriptsize\ttfamily\color{white},
  backgroundcolor=\color{lightgray},
  commentstyle = \color{darkgray},
  keywordstyle = \color{white},
  rulecolor    = \color{black},
  stringstyle  = \color{white}
  commentstyle=\color{lightgray}\ttfamily
}

\lstset{
   language=bash,
   extendedchars=true,
   basicstyle=\scriptsize\ttfamily,
   showstringspaces=false,
   showspaces=false,
   tabsize=2,
   showtabs=false,
   captionpos=b,
}

\usepackage{lipsum} 

\title{GEOMETRIC REASONING IN HUMAN TOOL USE}
\shorttitle{}
\author{
  \small{by}\\
  Adrian Ionita (ID: 1057404)
}
\affiliation{Supervisor: Dietmar Heinke\\
\vfill
\includegraphics[width=40mm]{figures/bham_logo.png}
\vfill
Submitted in partial fulfilment\\
of the requirements for the degree of\\ 
Master of Science
\vfill
School of Psychology\\
University of Birmingham\\
Birmingham, UK\\
September 2016
}
\abstract{
Tool use is a fundamental human ability that distinguishes man from other species.
This project investigates a potential building block for a computational model of tool use reasoning.
The building block focuses on solving geometric constraints for interactions between tools and recipient objects. 

Two approaches are considered:
\begin{enumerate}
\itemsep0em
\item An exhaustive approach using a physics engine to find compatible configurations of tool-object interaction by searching through of all possible alternatives. 
\item A heuristic model which reduces the search space by using visual cues of tool and object parts.
\end{enumerate}

The heuristic method employs a novel surface similarity technique which is able to assess parts having complementary geometries. 
Both exhaustive and heuristic models have potential for simulating human behaviour, but require further development and comparison with human behaviour.
}

\begin{document}
\pagenumbering{gobble}
\thispagestyle{otherpage}
\maketitle

\section{\normalfont Acknowledgements}
I would like to thank Dietmar Heinke, project supervisor, for his guidance, support and making tool-use a very interesting topic to work on. 
Thanks also go to Fran\c{c}ois Osiurak, the author of the four constraints theory, for his availability to clarifying concepts from his work. 
Last but not least, I would thank Rhiannon Davies for being the first subjected to proof-reading my work and providing constructive criticism.
\clearpage

\pagenumbering{arabic}

\tableofcontents
\clearpage
\listoffigures
\clearpage

\include*{ExhaustiveSearch}
\shorttitle{Exhaustive Search}
\clearpage
\include*{HeuristicSearch}

\clearpage
\section{Conclusion}
\shorttitle{Conclusion}
This project has developed two approaches for a computational model of geometric reasoning in human tool use.
The models would solve a task involving a pair of novel objects in which one object is used to lift and place the other at a new location. 
The task requires solving geometric constraints for tool-object interaction, and is best described as a process of technical reasoning.

In an exhaustive search model, a physics engine was used to find valid tool-object interaction by searching through all possible configurations.  
Although the search returns valid solutions for the lifting task, the model proves computationally unfeasible and unrealistic of human ability. 
Apart from long execution times the model would return solutions that are not physically possible to reach. 

Heuristic approaches were considered for a second model based on visual cues of object and tool parts.  
The underlying hypothesis was that parts which fit would share similar geometries. 
A novel surface similarity technique was presented to assess object parts from 3D data points.
Due to project time constraints, the approach was not fully implemented into a model demonstrating lifting solutions.
Only the visual search step of the process was implemented, but requires further investigation as it employs a simplified algorithm which produces false positive matches. Visual search results should also be formally verified against observed human behaviour.

Both the exhaustive search and heuristic models can be improved to provide better results.
The exhaustive search's simulation time can be reduced by running code on multiple machines and aggregating results.
The heuristic model can be improved through better analysis of objects geometries, but also by considering forces generated by user actions (i.e. lifting).
Nevertheless, even with the above improvements, model results require formal comparison to observed human behaviour performing similar tool-use tasks.  
\clearpage

\include*{Appendix}
\shorttitle{}
\printbibliography
\end{document}

